% Claim proof file for row claim_3_5 -- "BifurcationAlternative".
% Auto-generated stub by scaffold/solving_current_row.py at row start. A
% manager agent dedicated to writing the TeX proof fills in the body below.
% The proof must:
%   - Stay close to the lecture notes' proof if one exists (always search
%     the LN first for a `\begin{proof}` block immediately following the
%     claim; copy / lightly edit when present).
%   - Otherwise use the LN paradigm and cite prior claims by their `ref`.
%   - Be self-contained enough that a mathematician can follow it without
%     re-reading the LN.
%   - Have no `sorry`, `omitted`, or "obvious" shortcuts.
%
% Compiles standalone via the `subfiles` package.

\documentclass[main]{subfiles}

\begin{document}
% ref: claim_3_5 (proof, with statement restated for readability)
% title: BifurcationAlternative
\def\rowref{claim\_3\_5}
\def\rowtitle{BifurcationAlternative}
\phantomsection\label{claim_3_5}
%
% Cross-references: see claim_<ref>_statement_<title>.tex in this folder
% for the corresponding statement.

% --- statement (restated) ---------------------------------------------------
\begin{Prp}\label{prp:bifurcations_alternative}
  Let $G = \lt J, V, E, L \rt$ be a CDMG. For $v, w, c \in V \cup J$:
  there exists a bifurcation between $v$ and $w$ in $G$ with source $c$ if and only if $v \ne w$ and $c \in \Anc^{G_{\doit(w)}}(v) \sm \{v\}$ and $c \in \Anc^{G_{\doit(v)}}(w) \sm \{w\}$.
\end{Prp}

% --- proof ------------------------------------------------------------------
\begin{proof}
% TODO: write the proof body.
\end{proof}

\end{document}
